\documentclass[11pt]{article}
\usepackage[margin=1in]{geometry}
\usepackage{amsmath,amssymb,amsthm}
\usepackage{hyperref}
\newtheorem{theorem}{Theorem}[section]
\newtheorem{lemma}[theorem]{Lemma}
\newtheorem{proposition}[theorem]{Proposition}
\title{An affine Singer layer on $\mathbb{F}_2^6$, its $(T,D,K)$ triple, and $\langle T,D\rangle = \mathrm{AGL}(1,64)$}
\author{Baohua Li}
\date{\today}
\begin{document}
\maketitle
\begin{abstract}
Let $V=\mathbb{F}_2^6$ carry a fixed ordered basis $G=(g_0,\dots,g_5)$, and define the
affine map $T(x)=63\oplus \mathrm{coords}_G(x)$, where
$\mathrm{coords}_G(x)$ is the coordinate vector of $x$ in $G$, read as a six-bit word with the
first coordinate most significant. We show that $T$ has one fixed point $38$ and one $63$-cycle
(its linear part is a Singer cycle with characteristic polynomial
$x^6+x^5+x^3+x^2+1$); that $D=(\mathrm{id}\oplus T)^{-1}$ and
$K=(\mathrm{id}\oplus D)^{-1}$ are defined with cycle types $D=1+63$, $K=1+7\cdot 9$,
fixed points $\{38\},\{41\},\{63\}$, and the third-iterate identity
$(\mathrm{id}\oplus K)^{-1}=T$; and that $\langle T,D\rangle=\mathrm{AGL}(1,64)$,
whose order is $4032=64\cdot 63$. The analogous construction on $2k$-bit words is treated in \S\ref{sec:2k}: the cycle-type
existence for every $k$ is reduced to a prescribed-trace primitive-element (PTE) theorem, whose self-contained proof is given in Appendix~\ref{app:pte}
(historically, the even-characteristic case is covered by \cite{Mor} and
\cite[Theorem~1, p.~1]{Coh90}, and it also follows from an equivalent square-order formulation in
\cite[Theorem~1.5, p.~170]{Coh}); the generation
$\langle T,D\rangle=\mathrm{AGL}(1,2^{2k})$ holds for every $k$, and the concrete coordinate presentation is verified for $k=1,2,3,4$ (orders $12,240,4032,65280$). The linear
part of $K$ has characteristic polynomial an irreducible degree-$2k$ factor of the
cyclotomic polynomial $\Phi_{2^k+1}$. The main case is elementary over
$\mathbb{F}_{64}$; the even-characteristic existence input is proved self-containedly in Appendix~\ref{app:pte}, with historical references given for attribution.
\end{abstract}

\section{Introduction}\label{sec:intro}
The $64$ words of $\mathbb{F}_2^6$ support a classical combinatorial structure with
operations of complement, reversal, and a recombination that reads a word through its
coordinates in a fixed generator basis. This note isolates one such operation --- the
element transformation $T$ --- and shows that it is exactly an affine Singer map (layer) --- that is, an affine map whose
linear part is a Singer cycle. The
associated triple $(T,D,K)$ and the group $\langle T,D\rangle$ then reduce to
elementary finite-field statements, with a uniform analogue for $2k$-bit systems
(an existence theorem for all $k$, \S\ref{sec:2k} and Appendix~\ref{app:pte}; the concrete coordinate presentation is verified for $k=1,2,3,4$). Related families of
permutations over $\mathbb{F}_2^n$ with prescribed cycle structure and algebraic
degree were studied by Gravel, Panario and Thomson \cite{GPT}. The underlying
even-characteristic existence statement is classical (see \cite{Mor} and
\cite[Theorem~1, p.~1]{Coh90}; an equivalent square-order formulation is \cite[Theorem~1.5, p.~170]{Coh});
Appendix~\ref{app:pte} gives a self-contained proof.
Apart from these historical attributions, the paper is self-contained; the classical
names of the words are not used.

\section{Notation}\label{sec:not}
$V=\mathbb{F}_2^6$; coordinates $x=(x_1,\dots,x_6)$ with $x_1$ the most significant
bit (value $32$) and $x_6$ the least (value $1$); a word is identified with its
integer value in $\{0,\dots,63\}$. Basis $G=(\texttt{63},\texttt{33},\texttt{18},
\texttt{49},\texttt{1},\texttt{21})$ (a fixed ordered basis --- the six words are linearly independent over $\mathbb{F}_2$, by Gaussian elimination; non-uniqueness is discussed in \S\ref{sec:2k}); write $g_j$ for
the $j$-th entry, $j=0,\dots,5$.
$\mathrm{coords}_G(x)=(c_1,\dots,c_6)$ is the unique vector with
$x=\bigoplus_{i=1}^6 c_i\,g_{i-1}$, read as the six-bit word $c_1\cdots c_6$
($c_1$ most significant); $\oplus$ is bitwise XOR, i.e.~addition in $\mathbb{F}_2^n$.

\section[Affine Singer map]{$T$ is an affine Singer map}\label{sec:T}
\subsection*{3.1 Definition and basic checks}
$T(x)=63\oplus \mathrm{coords}_G(x)$. Verified identities:
$T(0)=63$, $T(41)=0$, $T(38)=38$; for example
$\mathrm{coords}_G(38)=(0,1,1,0,0,1)$, read as $011001=25$, so $T(38)=63\oplus25=38$,
and $\mathrm{coords}_G(41)=(1,1,1,1,1,1)=63$, so $T(41)=63\oplus63=0$.
\begin{lemma}\label{lem:aff}
$T$ is affine: $T(x)=L_Tx+b$ with $b=T(0)=63$ and $L_T(x)=T(x)\oplus 63$.
\end{lemma}
\begin{lemma}\label{lem:fix}
$\mathrm{Fix}(T)=\{38\}$; on the remaining $63$ points $T$ is a single cycle
(type $1+63$, order $63$), equivalently $L_T$ has order $63$ (Proposition~\ref{prop:singer}).
\end{lemma}
\begin{proposition}\label{prop:singer}
Identify $(V,+)$ with $\mathbb{F}_{64}$; then $L_T$ is
multiplication by a primitive element of $\mathbb{F}_{64}$, a Singer cycle \cite{LN}. The
characteristic polynomial of $L_T$ is
\begin{equation}\label{eq:pT}
p_T(t)=t^6+t^5+t^3+t^2+1,
\end{equation}
computed directly as $\det(tI+L_T)$ (code appendix), primitive of degree $6$.
\end{proposition}

\subsection*{3.2 Powers and $63=7\cdot 9=3\cdot 21$}
Since $L_T$ has order $63$, $T^d$ for $d\mid 63$ has type $1+d\times(63/d)$ (one fixed point plus $d$ cycles of length $63/d$):
$T^1=1+63$, $T^3=1+3\cdot21$, $T^7=1+7\cdot9$, $T^9=1+9\cdot7$,
$T^{21}=1+21\cdot3$ (verified).

\section[The (T,D,K) triple and the fixed-point operator]{The $(T,D,K)$ triple and the fixed-point operator}\label{sec:tdk}
For a permutation $f$ of $V$ such that $\mathrm{id}\oplus f$ is again a permutation,
define the \emph{fixed-point operator} $F(f)=(\mathrm{id}\oplus f)^{-1}$, where
$(\mathrm{id}\oplus f)(x)=x\oplus f(x)$. By Proposition~\ref{prop:singer},
identify the linear part of $T$ with multiplication by a primitive
$\alpha\ne1$. Then $\mathrm{id}\oplus T$ has nonzero slope
$1\oplus\alpha$ and is therefore a permutation. Its inverse $D=F(T)$ has slope
$(1\oplus\alpha)^{-1}$, and
$1\oplus(1\oplus\alpha)^{-1}=\alpha/(1\oplus\alpha)\ne0$, so
$\mathrm{id}\oplus D$ is also a permutation. Thus
$D=F(T)=(\mathrm{id}\oplus T)^{-1}$ and $K=F(D)=(\mathrm{id}\oplus D)^{-1}$ are
defined; since $F^3(T)=T$ below, $\{T,D,K\}$ is the $F$-orbit of $T$.
Verified facts:
\begin{itemize}
\item cycle types $T=D=1+63$, $K=1+7\cdot 9$ (seven $9$-cycles plus one fixed point);
\item fixed points $\mathrm{Fix}(T)=\{38\}$, $\mathrm{Fix}(D)=\{41\}$,
$\mathrm{Fix}(K)=\{63\}$;
\item $F(K)=T$, i.e.\ $(\mathrm{id}\oplus K)^{-1}=T$; hence $F^3=\mathrm{id}$ with
$T\to D\to K\to T$;
\item offset-fixed three-cycle: $D(0)=38=\mathrm{Fix}(T)$,
$K(0)=41=\mathrm{Fix}(D)$, $T(0)=63=\mathrm{Fix}(K)$;
\item linear parts: since $L_D=(I+L_T)^{-1}$, its characteristic polynomial is
the reciprocal of \eqref{eq:pT}, namely
$p_D(t)=t^6p_T(1/t)=t^6+t^4+t^3+t+1$ (primitive). Separately,
$p_T(t+1)=p_T(t)$ in characteristic two (directly,
$(t+1)^6+(t+1)^5+(t+1)^3+(t+1)^2+1$ reduces to
$t^6+t^5+t^3+t^2+1$); $L_K$ has $\Phi_9(t)=t^6+t^3+1$
(cyclotomic, non-primitive, order $9$).
\end{itemize}
\begin{lemma}\label{lem:fp}
For $g\in\{T,D,K\}$, the unique fixed point of $g$ is $F(g)(0)$, i.e.~
$\mathrm{Fix}(g)=\{F(g)(0)\}$.
\end{lemma}

\begin{lemma}\label{lem:F3}
In characteristic two, the fractional linear map $z\mapsto1/(1+z)$ has order three:
$z\mapsto1/(1+z)\mapsto1+1/z\mapsto z$ (where defined), and for the operator $F(f)=(\mathrm{id}\oplus f)^{-1}$ on affine maps
$f(x)=\alpha x\oplus c$ with $\alpha,1\oplus\alpha\ne0$, one has $F^3(f)=f$; the
slopes cycle as $\alpha\mapsto(1\oplus\alpha)^{-1}\mapsto1\oplus\alpha^{-1}\mapsto\alpha$.
\end{lemma}
\begin{proof}
Set $\beta=(1\oplus\alpha)^{-1}$ and $\eta=1\oplus\alpha^{-1}$. Then
$F(f)(x)=\beta x\oplus\beta c$,
$F^2(f)(x)=\eta x\oplus\alpha^{-1}c$, and
$F^3(f)(x)=\alpha x\oplus c=f(x)$. In particular,
$\alpha\beta\eta=1$.
\end{proof}

\textbf{Remark.} The fractional-linear statement holds in characteristic two and fails in odd
characteristic; the affine identity $F^3=\mathrm{id}$ is genuinely characteristic-two: over a field of odd
characteristic, whenever the iterates are defined (equivalently,
$(1+\alpha)(2+\alpha)(3+2\alpha)\neq0$),
$F^3(f)(x)=\frac{\alpha+2}{2\alpha+3}x-\frac{c}{2\alpha+3}\neq\alpha x+c=f(x)$
whenever $c\neq0$. Thus the
$\mathbb{F}_2$ setting is essential for the third-order affine cycle, not merely a
convenience.

The identity is affine-general: by Lemma~\ref{lem:F3}, in field coordinates
$T(x)=\alpha x\oplus c$, so
$F(T)(0)=D(0)=(1\oplus\alpha)^{-1}c$ solves $x=\alpha x\oplus c$, i.e.\
$D(0)=\mathrm{Fix}(T)$, and the slope cycle
$\alpha\mapsto(1\oplus\alpha)^{-1}\mapsto1\oplus\alpha^{-1}\mapsto\alpha$ gives the
two remaining identities. For the six-bit layer the values are $F(T)(0)=38$,
$F(D)(0)=41$, $F(K)(0)=63$.
Conceptually, $T$ is the launching map --- an affine Singer layer read through the
fixed generator basis (a change of basis) --- while the fixed-point operator $F$ keeps
the third-order cycle running: $T\mapsto D\mapsto K\mapsto T$, carrying the fixed
point $38\mapsto41\mapsto63\mapsto38$.

\section[The group generated by T and D]{The group $\langle T,D\rangle$}\label{sec:group}
\begin{theorem}\label{thm:agld}
$\langle T,D\rangle=\mathrm{AGL}(1,64)$; its order is $4032=64\cdot63$.
\end{theorem}
\begin{proof}
By Proposition~\ref{prop:singer}, choose the field identification in which the linear
part of $T$ is multiplication by a primitive $\alpha$ (order $63$); then
$T(x)=\alpha x\oplus c$ with $c=T(0)=63\ne0$. Since
$(\mathrm{id}\oplus T)(x)=(1\oplus\alpha)x\oplus c$, the inverse is affine,
$D(x)=\beta x\oplus\beta c$ with $\beta=(1\oplus\alpha)^{-1}$. Every word of
$\langle T,D\rangle$ is therefore affine with multiplier in
$\mathbb{F}_{64}^*=\langle\alpha\rangle$, so
$\langle T,D\rangle\subseteq\mathrm{AGL}(1,64)$.
For the converse, $T\circ D$ and $D\circ T$ have the same multiplier $\alpha\beta$:
$T\circ D(x)=\alpha\beta x\oplus\alpha\beta c\oplus c$ and $D\circ T(x)=\alpha\beta x$,
so $T\circ D(x)=D\circ T(x)\oplus c\beta$ (as $\alpha\beta\oplus1=\beta$), with
$c\beta=c/(1\oplus\alpha)\ne0$;
hence the commutator $[T,D]=TDT^{-1}D^{-1}$ has multiplier $1$ and is the nonzero translation
$x\mapsto x\oplus c/(1\oplus\alpha)$. Conjugating it by $T^j$,
$j=0,\dots,62$, gives the translations $x\mapsto x\oplus\alpha^j c/(1\oplus\alpha)$,
which, as $\alpha$ is primitive, are all nonzero translations together with the
identity. Since the powers of $T$ supply every multiplier
$\mathbb{F}_{64}^*=\langle\alpha\rangle$, together with these translations
$\langle T,D\rangle$ contains
$\{x\mapsto ax\oplus b: a\in\mathbb{F}_{64}^*,\,b\in\mathbb{F}_{64}\}
=\mathrm{AGL}(1,64)$ (cf.\ \cite{DM}).
\end{proof}

\section[Uniform statement for 2k-bit words]{Uniform statement for $2k$-bit words}\label{sec:2k}
Let $V_{2k}=\mathbb{F}_{2^{2k}}$; write $M=2^{2k}-1=(2^k-1)(2^k+1)$. In field
coordinates an affine Singer layer is a map $T(x)=\alpha x\oplus b$ whose linear part
is multiplication by a primitive element $\alpha$ (order $M$).
\begin{lemma}\label{lem:ord}
For every $k\ge1$, $\operatorname{ord}_{2^k+1}(2)=2k$.
\end{lemma}
\begin{proof}
Let $m=\operatorname{ord}_{2^k+1}(2)$. Since $2^k\equiv-1\pmod{2^k+1}$,
we have $m\mid2k$, while $m\nmid k$. Suppose, for contradiction, that
$m<2k$. Since $m\mid2k$ but $m\nmid k$, $m$ is even; write $m=2d$
with $d\mid k$, and $m<2k$ gives $d<k$. The subgroup
$\langle2\rangle\subset(\mathbb Z/(2^k+1)\mathbb Z)^\times$ is cyclic of order
$2d$, so its unique element of order $2$ is $2^d$. But $2^k\equiv-1$ also has
order $2$ in this subgroup; hence $2^d\equiv-1\pmod{2^k+1}$. Thus
$2^k+1\mid 2^d+1$, contradicting
$0<2^d+1<2^k+1$. Therefore $m=2k$.
\end{proof}
For \(q=2^k\), define
\[
\operatorname{PTE}(k):\quad
\exists\beta\in\mathbb{F}_{q^2}^*,\quad
\operatorname{ord}(\beta)=q^2-1,\quad
\operatorname{Tr}_{\mathbb{F}_{q^2}/\mathbb{F}_q}(\beta)=1.
\]
The following theorem supplies the required existence input; its proof is given in Appendix~\ref{app:pte}.
\begin{theorem}[PTE existence]\label{thm:pte}
For every \(k\ge1\), \(\operatorname{PTE}(k)\) holds.
\end{theorem}
\begin{proof}
See Appendix~\ref{app:pte}.
\end{proof}
\begin{theorem}\label{thm:2k}
For every $k\ge1$ there is an affine Singer layer $T$ (with nonzero translation) on
$\mathbb{F}_{2^{2k}}$ for which $D=(\mathrm{id}\oplus T)^{-1}$ and $K=(\mathrm{id}\oplus D)^{-1}$ are defined,
$T$ and $D$ have cycle type $1+M$, $K$ has cycle type
$1+(2^k-1)\times(2^k+1)$, and the fixed-point operator satisfies $F^3=\mathrm{id}$
with $T\mapsto D\mapsto K\mapsto T$.
\end{theorem}
\begin{proof}[Proof (summary)]
Take $T(x)=\alpha x\oplus b$ with $b\ne0$. The linear parts of $T$, $D$, $K$ are
$\alpha$, $(1\oplus\alpha)^{-1}$, and
$\kappa=1\oplus\alpha^{-1}=(1\oplus\alpha)/\alpha$
respectively (independent of the translation $b$); the three slopes are cycled by
the map $z\mapsto 1/(z\oplus1)$, so $F^3=\mathrm{id}$. Since
$\mathbb{F}_{2^{2k}}^*$ is cyclic, an affine map with slope $\mu\ne1$ has cycle type
$1+\frac{M}{\mathrm{ord}(\mu)}\times\mathrm{ord}(\mu)$ (the translation only moves the
fixed point; a shorter orbit would force $T^d$, $d<\mathrm{ord}(\mu)$, to have a second
fixed point, impossible since $\mu^d\ne1$); the stated types therefore follow once
$\mathrm{ord}(\alpha)=\mathrm{ord}(1\oplus\alpha)=M$ and
$\mathrm{ord}(\kappa)=2^k+1$. Existence of such an $\alpha$ follows from Theorem~\ref{thm:pte}. Choose
$\beta\in\mathbb{F}_{q^2}^*$ with
$\mathrm{ord}(\beta)=M$ and $\operatorname{Tr}_{\mathbb{F}_{q^2}/\mathbb{F}_q}(\beta)=1$,
and put $\gamma=\beta^{q-1}$. Then $\gamma\in\mathbb{F}_q$,
$\mathrm{ord}(\gamma)=q+1$, and $1\oplus\gamma=\beta^{-1}$; hence
\[
\tau=\gamma\oplus\gamma^{-1}
=(1\oplus\gamma)^q(1\oplus\gamma)
=\beta^{-(q+1)}
\]
has order $q-1$, so $\tau$ is primitive in $\mathbb{F}_q$. The $\tau$-collapse
now applies with $\alpha=(1\oplus\gamma)^{-1}$.
\end{proof}
Since the $K$-slope has order $2^k+1$, the degree over $\mathbb{F}_2$ of its
minimal polynomial is $\operatorname{ord}_{2^k+1}(2)=2k$ by
Lemma~\ref{lem:ord}; this polynomial is therefore an irreducible degree-$2k$
factor of $\Phi_{2^k+1}$.
\textbf{Remark (the $\tau$-collapse).}
For an element $\gamma$ of order $2^k+1$, the Frobenius identities
$(1\oplus\gamma)^{2^k-1}=\gamma^{-1}$ and $(1\oplus\gamma)^{2^k+1}=\tau$,
together with $\gcd(2^k-1,2^k+1)=1$, give
$\mathrm{ord}(1\oplus\gamma)=\mathrm{lcm}(\mathrm{ord}(\gamma^{-1}),\mathrm{ord}(\tau))
=\mathrm{lcm}(2^k+1,2^k-1)=M$ whenever $\tau$ is primitive in $\mathbb{F}_{2^k}$;
the pair $(u^{2^k-1},u^{2^k+1})=(\gamma^{-1},\tau^{-1})$ for
$u=\gamma/(1\oplus\gamma)$ gives the same order for $u$. Since
$\alpha=(1\oplus\gamma)^{-1}$ and
$\kappa=1\oplus\alpha^{-1}=\gamma$, once $\gamma$
has order $2^k+1$ the three order conditions on $\alpha$ reduce to the single
condition that $\tau$ be primitive. The order assumption on $\gamma$ is
essential; it is not implied by $\tau$ being primitive alone.
Existence of such a $\gamma$ is the classical square-order statement
(see \cite{Mor} and \cite[Theorem~1, p.~1]{Coh90}; for an equivalent square-order formulation, see
\cite[Theorem~1.5, p.~170]{Coh}); it is also supplied directly by
Theorem~\ref{thm:pte}.
For $q=2^k$ the appendix proof is self-contained, so the paper does not use these
classical results as external existence inputs.

The theorem is an existence statement in admissible coordinates (equivalently, for an
admissible ordered generator basis). The argument of \S\ref{sec:group} applies verbatim
with $q=2^{2k}$ (taking $b\ne0$ in the proof), so the generation
$\langle T,D\rangle=\mathrm{AGL}(1,2^{2k})$ holds for every $k$. For the concrete
coordinate presentation of \S\ref{sec:T} the orders and cycle data were verified
for $k=1,2,3,4$ (the bases are listed in the appendix):
$k=1$ (2 bits) order $12=4\cdot3$, $K$ type $1+1\cdot3$, factor
$\Phi_3$ (basis $(1,3)$: $\mathrm{Fix}(T)=\{1\}$, $\mathrm{Fix}(D)=\{2\}$,
$\mathrm{Fix}(K)=\{3\}$); $k=2$ (4 bits) order $240=16\cdot15$, $K$ type
$1+3\cdot5$, factor $\Phi_5$; $k=3$ (6 bits) order $4032=64\cdot63$, $K$ type
$1+7\cdot9$, factor $\Phi_9$ (this paper); $k=4$ (8 bits) order
$65280=256\cdot255$, $K$ type $1+15\cdot17$, linear part an irreducible degree-$8$
factor of $\Phi_{17}$.

The basis is not canonical even in dimension $6$: among the $720$ orderings of these
six generators, $572$ do not even make $\mathrm{id}\oplus T$ a permutation; of the $148$
that do, only $54$ give $T$ of order $63$, and only $9$ give the exact cycle-type triple
$T=D=1+63$, $K=1+7\cdot9$. For instance $(4\,6)$ keeps $T=1+63$ but changes
$\mathrm{Fix}(T)$ from $38$ to $11$ and the cycle types of $D$ and $K$ to $D=1+3\cdot21$,
$K=1+63$. Hence the particular entries of $G$ and the anchors $\{38,41,63\}$ are coordinates
of one representative, not invariants, and no uniqueness statement is made. The structural results of \S\ref{sec:group}--\ref{sec:2k},
however, are representative-independent: $\langle T,D\rangle=\mathrm{AGL}(1,64)$, the
third-iterate identity $F^3=\mathrm{id}$, and the $2k$-bit existence statement hold for any
admissible basis; only the numerical anchors change with the basis.

\begingroup
\emergencystretch=1em
\noindent\textbf{Remark (constructiveness).}\\
Theorem~\ref{thm:2k} is an existence statement; its existence input is proved
in Appendix~\ref{app:pte}. The $\tau$-collapse turns it into a finite search with
certification. For $\tau=\gamma\oplus\gamma^{-1}\in\mathbb{F}_{2^k}$, the
Artin--Schreier identity
(\emph{$x^2+\tau x+1$ is irreducible over $\mathbb{F}_{2^k}$ iff
$\mathrm{Tr}_{\mathbb{F}_{2^k}/\mathbb{F}_2}(1/\tau)=1$}) gives $2^{k-1}$
candidate values of $\tau$. For a candidate, fix an explicit representation of
$\mathbb{F}_{2^{2k}}$ and take a root $\gamma$ of $x^2+\tau x+1$
(found, for instance, by exhaustive search); the collapse gives the required orders for $1\oplus\gamma$ and
$\gamma/(1\oplus\gamma)$ precisely when $\tau$ is primitive \emph{and}
$\mathrm{ord}(\gamma)=2^k+1$. The second condition is not automatic from the
first: for example, when $k=6$ and
$\mathbb{F}_{64}=\mathbb{F}_2[t]/(t^6+t+1)$, the element
$\tau=t^5+t^3+1$ is primitive and $\mathrm{Tr}(1/\tau)=1$, but the roots of
$x^2+\tau x+1$ have order $13$, not $65$. The theorem of
Appendix~\ref{app:pte} guarantees that a jointly successful candidate exists, so a
deterministic enumeration that tests both conditions terminates. Certifying them requires factoring $2^k-1$ and
$2^k+1$ (or an equivalent exact-order test). Only a factorization-free closed
form (the classical explicit-primitive-element problem) is left open.
A necessary restatement is that $u=1/\tau$ is a primitive element of
$\mathbb{F}_{2^k}$ with $\mathrm{Tr}_{\mathbb{F}_{2^k}/\mathbb{F}_2}(u)=1$;
equivalently, its degree-$k$ minimal polynomial has $x^{k-1}$-coefficient
$1$. This condition alone is still not sufficient for admissibility, because
the root-order condition above must also be checked.
\endgroup

\section{Remarks}\label{sec:rem}
The fixed point $38$ is a ``power-layer'' origin (domain zero after relabeling), while
the value $0$ is the additive origin; no interpretive claim is made here. In the value set $\{0,\dots,63\}$ the $9$-multiples are exactly the $8$ self-doubled words ($x_i=x_{i+3}$), a $3$-dimensional subspace; the $63$-cycle of $T$ is parametrized by the cyclic group $\mathbb{Z}/63\cong\mathbb{Z}/7\times\mathbb{Z}/9$ (CRT decomposition); the all-ones word $63$ represents the zero class modulo $63$, hence the $(0,0)$ corner of this decomposition. The
cyclotomic factor $\Phi_{2^k+1}$ mirrors the identity $(2^k-1)(2^k+1)=2^{2k}-1$; the
single fixed point of $K$ accounts for the remaining point, so $K$ permutes all
$2^{2k}$ words. Two questions are left open: the orbit structure of the fixed-point operator $F$
(a bijection of order three on the admissible affine maps in characteristic two)
has not been classified, and whether analogous cycle-type/factorization
correspondences hold for other factorizations of $q^n-1$ is not pursued here.

\begin{thebibliography}{9}
\bibitem{LN} R.~Lidl, H.~Niederreiter, \emph{Finite Fields}, 2nd ed., Cambridge, 1997.
\bibitem{DM} J.~D.~Dixon, B.~Mortimer, \emph{Permutation Groups}, Springer, 1996.
\bibitem{Coh} S.~D.~Cohen, \emph{The orders of related elements of a finite field}, Ramanujan J.~7 (2003) 169--183, doi:10.1023/A:1026295128600.
\bibitem{Mor} O.~Moreno, \emph{On the existence of a primitive quadratic of trace $1$ over $GF(p^m)$}, J.~Combin.~Theory Ser.~A 51 (1989) 104--110.
\bibitem{Coh90} S.~D.~Cohen, \emph{Primitive elements and polynomials with arbitrary trace}, Discrete Math.~83 (1990) 1--7.
\bibitem{GPT} C.~Gravel, D.~Panario, D.~Thomson, \emph{Unicyclic strong permutations}, arXiv:1809.03551 (2018); Cryptogr.~Commun.~11 (2019) 1211--1231.
\end{thebibliography}

\appendix
\section{Computational appendix}\label{app:code}
All numerical claims (fixed points, cycle types, group order $4032$, and the
characteristic polynomials of $L_T$, $L_D$, and $L_K$) were verified
computationally; the verification code is available from the author, and the finite k<=29 verifier is included as \texttt{verify\_k1\_29.py}. The characteristic
polynomials were computed as $\det(tI+L)$ via evaluation over $\mathbb{F}_8$ and
interpolation. Relative to the standard basis $e_1,\dots,e_6$, the linear part of
the main map is
\[
L_T=\begin{pmatrix}
0&0&1&0&0&0\\
1&1&0&1&1&0\\
0&0&1&0&1&0\\
0&1&1&1&1&0\\
1&0&1&1&0&1\\
0&0&1&1&0&0
\end{pmatrix}.
\]
\emph{Formalization.} The affine group-generation argument (for a field $\mathbb K$ with a primitive
element $\alpha$ and $c\ne0$, the maps $T$ and $D$ generate $\mathrm{AGL}(1,\mathbb K)$, the core
of \S\ref{sec:group} and \S\ref{sec:2k}) and the algebraic lemmas used in it (the commutator
is a translation; conjugation of translations by $T$) have been formalized and
machine-checked in Lean~4 using mathlib. An additional conditional Lean file
formalizes the field-level Frobenius identities and the resulting
order-theoretic $\tau$-collapse under an explicit order interface; it does not
formalize the paper proof of the existence theorem. The concrete cycle-type data of
\S\ref{sec:T}--\ref{sec:tdk} were verified by independent computation. The existence
part of Theorem~\ref{thm:2k} is proved in Appendix~\ref{app:pte}; its analytic part
is not formalized in Lean, and its finite range is discharged by the deterministic
certificate there. Thus Cohen's theorem is not used as a proof dependency; the
classical references are cited only for historical attribution. The Lean source is
available from the author; the finite verifier is included as
\texttt{verify\_k1\_29.py}.

\subsection*{Bases for the coordinate presentation}
For the coordinate presentation $T(x)=(2^{2k}-1)\oplus\mathrm{coords}_G(x)$ on
$2k$ bits the following bases give the stated group order $\langle T,D\rangle$ and
$K$-cycle type:
\begin{center}
\begin{tabular}{c|c|c|c|c}
$k$ & basis $G$ & $\langle T,D\rangle$ & $K$ type & factor \\
\hline
1 & $(1,3)$ & $12=4\cdot3$ & $1+1\cdot3$ & $\Phi_3$ \\
2 & $(1,2,5,13)$ & $240=16\cdot15$ & $1+3\cdot5$ & $\Phi_5$ \\
3 & $(63,33,18,49,1,21)$ & $4032=64\cdot63$ & $1+7\cdot9$ & $\Phi_9$ \\
4 & $(134,136,241,71,34,65,3,143)$ & $65280=256\cdot255$ & $1+15\cdot17$ & a degree-$8$ factor of $\Phi_{17}$ \\
\end{tabular}
\end{center}
All four rows were verified by independent computation. The $720$-case count
was reproduced by enumerating all orderings of the six displayed vectors, solving
the corresponding $\mathbb{F}_2$ linear system for $\mathrm{coords}_G$, and
recording the cycle types of $T,D,K$ whenever the two inverses exist.

\subsection*{Explicit $\tau$-collapse data}
With $\mathbb{F}_{2^{2k}}=\mathbb{F}_2[t]/(f)$ and elements written as polynomials in $t$,
the following rows realize the order conditions of Theorem~\ref{thm:2k}:
\begin{center}
\begin{tabular}{c|c|c|c|c}
$k$ & $f$ & $\gamma$ & $\tau=\gamma\oplus\gamma^{-1}$ & $\alpha=(1\oplus\gamma)^{-1}$ \\
\hline
1 & $t^2+t+1$ & $t$ & $1$ & $t$ \\
2 & $t^4+t+1$ & $t^3$ & $t^2+t+1$ & $t$ \\
3 & $t^6+t+1$ & $t^2+t$ & $t^4+t^3+1$ & $t^5+t^3+t^2$ \\
4 & $t^8+t^4+t^3+t^2+1$ & $t^3+t^2+t+1$ & $t^7+t^4+t^3+1$ & $t^6+t^4+t^3+t^2+1$ \\
\end{tabular}
\end{center}
In each row $\mathrm{ord}(\gamma)=\mathrm{ord}(1\oplus\alpha^{-1})=2^k+1$, $\tau$ is
primitive in $\mathbb{F}_{2^k}$, and
$\mathrm{ord}(\alpha)=\mathrm{ord}(1\oplus\alpha)=2^{2k}-1$ (all verified). The defining polynomial $f$ is a field-representation choice and need not coincide with the characteristic polynomial $p_T$ of \eqref{eq:pT}; for instance, the $k=3$ row uses $f=t^6+t+1$, while $p_T(t)=t^6+t^5+t^3+t^2+1$.



\section{Self-contained proof for the even-characteristic case}\label{app:pte}
\begin{proof}[Proof of Theorem~\ref{thm:pte}]
The statement is classical (see \cite{Mor} and
\cite[Theorem~1, p.~1]{Coh90}; it also follows from the equivalent square-order formulation in
\cite[Theorem~1.5, p.~170]{Coh}); we give a self-contained proof in the
even-characteristic case used above. This is a reconstruction of known results (Moreno; Cohen), and no priority is claimed; it is included only to make the paper self-contained.
Put $q=2^k$, $L=\mathbb{F}_{q^2}$, $K_0=\mathbb{F}_q$, and
$M=q^2-1=(q-1)(q+1)$. Let
\[
E_1=\{x\in L:x^q+x=1\}.
\]
The trace map $L\to K_0$ is surjective and $K_0$-linear, so $|E_1|=q$. If
$x\in E_1$, then $x\notin K_0$ and its minimal polynomial over $K_0$ is
$X^2+X+xx^q$. Hence it is enough to prove that $E_1$ contains an element of
order $M$.

Let $\psi$ be the nontrivial additive character of $\mathbb F_2$, and put
$\psi_L(x)=\psi(\operatorname{Tr}_{L/\mathbb F_2}(x))$ and
$\psi_{K_0}(a)=\psi(\operatorname{Tr}_{K_0/\mathbb F_2}(a))$. For a multiplicative
character $\chi$ of $L^*$ define
\[
g_L(\chi)=\sum_{x\in L^*}\chi(x)\psi_L(x),\qquad
g_{K_0}(\chi)=\sum_{a\in K_0^*}\chi(a)\psi_{K_0}(a),\qquad
S(\chi)=\sum_{x\in E_1}\chi(x).
\]
If $\chi$ is nontrivial on $L^*$, then $|g_L(\chi)|=q$; if
$\chi|_{K_0^*}$ is nontrivial, then $|g_{K_0}(\chi)|=\sqrt q$. Indeed, the usual
squared-modulus calculation gives
\[
|g_L(\chi)|^2
=\sum_{t\in L^*}\chi(t)^{-1}\sum_{a\in L^*}\psi_L(a(1-t))=q^2,
\]
and the same argument in $K_0$ gives $q$.

We first compute $S(\chi)$. For $\alpha\in K_0$ set
$E_\alpha=\{x\in L:x^q+x=\alpha\}$; these sets partition $L$, and
$E_\alpha=\alpha E_1$ for $\alpha\ne0$. If $\chi^{q+1}\ne1$, then the
restriction of $\chi$ to $K_0^*$ is nontrivial; hence
\[
g_L(\chi)
=\sum_{\alpha\in K_0}\psi_{K_0}(\alpha)\sum_{x\in E_\alpha}\chi(x)
=S(\chi)g_{K_0}(\chi),
\]
so $|S(\chi)|=q/\sqrt q=\sqrt q$.

Suppose now that $\chi^{q+1}=1$ and $\chi\ne1$. Then $\chi|_{K_0^*}=1$.
Let $H=\{h\in L^*:h^{q+1}=1\}$; since $\gcd(q-1,q+1)=1$, the group $L^*$
is the direct product $K_0^*\times H$. Write $x=ch$ with $c\in K_0^*$ and
$h\in H$. The condition $x^q+x=1$ becomes
$c(h+h^{-1})=1$; thus $h\ne1$ and $c=(h+h^{-1})^{-1}$. This gives a
bijection between $H\setminus\{1\}$ and $E_1$, and therefore
\[
S(\chi)=\sum_{h\in H,\ h\ne1}\chi(h)
=\sum_{h\in H}\chi(h)-1=-1,
\]
because $\chi|_H$ is nontrivial.

We shall use the primitive-element indicator
\begin{equation}\label{eq:pte-indicator}
\mathbf 1_{\operatorname{ord}(x)=M}
=\frac{\varphi(M)}M
\sum_{\delta\mid M}\frac{\mu(\delta)}{\varphi(\delta)}
\sum_{\operatorname{ord}(\chi)=\delta}\chi(x).
\end{equation}
To see this, fix a generator $g$ of $L^*$ and write $x=g^a$. The left side
is the indicator of $\gcd(a,M)=1$. Its Fourier coefficient at a character
of order $\delta$ is the Ramanujan sum
$c_M(j)=\mu(\delta)\varphi(M)/\varphi(\delta)$; Fourier inversion gives
\eqref{eq:pte-indicator}. (The displayed identity for $c_M(j)$ follows by
M\"obius inversion.)

Assume for contradiction that $E_1$ has no element of order $M$. Summing
\eqref{eq:pte-indicator} over $E_1$ gives
\[
0=\frac{\varphi(M)}M
\sum_{\delta\mid M}\frac{\mu(\delta)}{\varphi(\delta)}S_\delta,
\qquad
S_\delta=\sum_{\operatorname{ord}(\chi)=\delta}S(\chi).
\]
Let $s=\omega(M)$ and $s'=\omega(q+1)$ be the numbers of distinct prime
divisors. The $\delta=1$ term is $q$. If $\delta\mid q+1$ and $\delta>1$,
then each character of order $\delta$ satisfies $\chi^{q+1}=1$, so the
corresponding contribution is $-\mu(\delta)$; summing over
$\delta\mid q+1$, $\delta>1$, gives $1$. For the remaining squarefree
divisors $\delta\nmid q+1$, the character sums are bounded by $\sqrt q$,
and their number is at most $2^s-2^{s'}$. Consequently,
\[
q+1\le(2^s-2^{s'})\sqrt q. \tag{A.2}
\]
Since $q+1>\sqrt q$, this implies
\[
2^s>2^{k/2}, \qquad\text{so}\qquad s>k/2.
\]
In particular $k<2s$. Let $p_1<p_2<\cdots$ be the primes and
$P_s=p_1\cdots p_s$. Since $M=2^{2k}-1$ has exactly $s$ distinct prime
divisors, $M\ge P_s$. We have
\[
P_{16}=32589158477190044730>2^{64}.
\]
For $n\ge16$, the ratio
\[
\frac{P_{n+1}/2^{4(n+1)}}{P_n/2^{4n}}
=\frac{p_{n+1}}{16}
\]
is greater than $1$, because $p_{n+1}\ge59$. Hence $P_s>2^{4s}$ for
every $s\ge16$. But $k<2s$ gives
\[
M=2^{2k}-1<2^{2k}<2^{4s},
\]
so $s\ge16$ is impossible. Thus $s\le15$ and $k<2s\le30$, i.e.
$k\le29$.

It remains to cover this finite range. In the table below, $m_k$ is an
irreducible polynomial of degree $k$ over $\mathbb F_2$, encoded by the
binary coefficient vector of the displayed hexadecimal integer; thus
$K_0=\mathbb F_2[x]/(m_k)$. The element $a_k\in K_0$ is encoded similarly.
For each row let $t$ be a root of $X^2+X+a_k$ in
$K_0[t]/(X^2+X+a_k)$. The deterministic verifier checks that $m_k$ is
irreducible, that $\operatorname{Tr}_{K_0/\mathbb F_2}(a_k)=1$, that
$t^{2^k}=t+1$, and that $t^M=1$ while $t^{M/p}\ne1$ for every prime
$p\mid M$. The last two checks show that $\operatorname{ord}(t)=M$; the
third shows that $\operatorname{Tr}_{L/K_0}(t)=1$. A verifier for the table, \texttt{verify\_k1\_29.py},
is included as an ancillary file; its SHA-256 hash is
\texttt{042CE7F74EDFF13ADE5C6285E42CB3D15491FC25}\allowbreak
\texttt{C65B99BC670D256047EE8D9F}.
\begin{center}
\small
\begin{tabular}{rrr}
$k$ & $m_k$ & $a_k$\\
\hline
1 & \texttt{0x3} & \texttt{0x1}\\
2 & \texttt{0x7} & \texttt{0x3}\\
3 & \texttt{0xb} & \texttt{0x5}\\
4 & \texttt{0x13} & \texttt{0xd}\\
5 & \texttt{0x25} & \texttt{0x13}\\
6 & \texttt{0x43} & \texttt{0x34}\\
7 & \texttt{0x83} & \texttt{0x47}\\
8 & \texttt{0x11b} & \texttt{0xc8}\\
9 & \texttt{0x203} & \texttt{0x2b}\\
10 & \texttt{0x409} & \texttt{0xcb}\\
11 & \texttt{0x805} & \texttt{0x525}\\
12 & \texttt{0x1009} & \texttt{0x39e}\\
13 & \texttt{0x201b} & \texttt{0x1ec4}\\
14 & \texttt{0x4021} & \texttt{0x3717}\\
15 & \texttt{0x8003} & \texttt{0x62bf}\\
16 & \texttt{0x1002b} & \texttt{0x1e6e}\\
17 & \texttt{0x20009} & \texttt{0x1155b}\\
18 & \texttt{0x40009} & \texttt{0x2ee84}\\
19 & \texttt{0x80027} & \texttt{0xb9af}\\
20 & \texttt{0x100009} & \texttt{0x6aee2}\\
21 & \texttt{0x200005} & \texttt{0x15750b}\\
22 & \texttt{0x400003} & \texttt{0x24a150}\\
23 & \texttt{0x800021} & \texttt{0x6c92f9}\\
24 & \texttt{0x100001b} & \texttt{0x389ec4}\\
25 & \texttt{0x2000009} & \texttt{0x6801ad}\\
26 & \texttt{0x400001b} & \texttt{0x2707d2e}\\
27 & \texttt{0x8000027} & \texttt{0x6e9ab1e}\\
28 & \texttt{0x10000003} & \texttt{0xfac7cd2}\\
29 & \texttt{0x20000005} & \texttt{0xc69b55e}\\
\end{tabular}
\end{center}
The table therefore supplies the required element for every $k\le29$.
Combining this with the contradiction for $k\ge30$ proves the theorem.
\end{proof}

\end{document}







